\section{Метод граничных элементов}
\label{sec:boundary-element-method}

\subsection{Основная идея метода граничных элементов}

\textbf{Метод граничных элементов (МГЭ)} --- численный метод решения краевых задач математической физики, в котором неизвестные аппроксимируются только на границе области, а не во всём её объёме.

Пусть требуется решить краевую задачу
$$
L u(\mathbf{x}) = f(\mathbf{x}), \qquad \mathbf{x}\in\Omega,
$$
с заданными граничными условиями на
$$
\Gamma = \partial\Omega.
$$

Основная идея МГЭ заключается в следующем:
\begin{enumerate}
	\item исходное дифференциальное уравнение преобразуется в \textbf{граничное интегральное уравнение};
	\item область $\Omega$ не разбивается на конечные элементы;
	\item разбивается только её граница $\Gamma$ на небольшие участки --- \textbf{граничные элементы};
	\item неизвестные на границе аппроксимируются простыми функциями;
	\item интегральное уравнение заменяется системой линейных алгебраических уравнений.
\end{enumerate}

Таким образом, для двумерной задачи вместо дискретизации двумерной области требуется дискретизация её одномерной границы, а для трёхмерной --- двумерной поверхности.

\subsection{Фундаментальное решение и функция Грина}

Ключевым элементом МГЭ является \textbf{фундаментальное решение} уравнения
$$
L G(\mathbf{x},\boldsymbol{\xi}) = -\delta(\mathbf{x}-\boldsymbol{\xi}),
$$
где $\delta$ --- дельта-функция Дирака, а $\boldsymbol{\xi}$ --- точка приложения единичного источника.

Физически $G(\mathbf{x},\boldsymbol{\xi})$ можно интерпретировать как реакцию системы в точке $\mathbf{x}$ на единичный точечный источник, расположенный в точке $\boldsymbol{\xi}$.

Для оператора Лапласа $\nabla^2$ фундаментальное решение имеет вид:
$$
G(\mathbf{x},\boldsymbol{\xi}) =
\begin{cases}
	-\dfrac{1}{2\pi}\ln r, & \text{в двумерном случае},\\[2mm]
	\dfrac{1}{4\pi r}, & \text{в трёхмерном случае},
\end{cases}
\qquad
r = |\mathbf{x}-\boldsymbol{\xi}|.
$$
Знак зависит от принятого определения оператора: при $L G = +\delta$ знаки меняются на противоположные.

Фундаментальное решение позволяет преобразовать исходное дифференциальное уравнение в интегральное уравнение, содержащее интегралы по границе области.

\subsection{Переход к граничному интегральному уравнению}

Рассмотрим для простоты задачу Дирихле для уравнения Лапласа:
$$
\nabla^2 u = 0 \quad \text{в } \Omega, \qquad
u = \bar u \quad \text{на } \Gamma.
$$

Используя вторую формулу Грина и фундаментальное решение, получаем \textbf{граничное интегральное уравнение} (формула Сохоцкого--Гёльдера):
$$
c(\boldsymbol{\xi})\,u(\boldsymbol{\xi})
+
\int_{\Gamma} u(\mathbf{x})\,
\frac{\partial G(\mathbf{x},\boldsymbol{\xi})}{\partial n}
\,d\Gamma
=
\int_{\Gamma} G(\mathbf{x},\boldsymbol{\xi})\,
\frac{\partial u(\mathbf{x})}{\partial n}
\,d\Gamma,
\qquad \boldsymbol{\xi}\in\Gamma.
$$

Здесь:
\begin{itemize}
	\item $u$ --- искомая функция;
	\item $\dfrac{\partial u}{\partial n}$ --- нормальная производная функции $u$;
	\item $G$ --- фундаментальное решение;
	\item $n$ --- внешняя нормаль к границе;
	\item $\boldsymbol{\xi}$ --- точка наблюдения;
	\item $c(\boldsymbol{\xi})$ --- коэффициент, зависящий от положения точки.
\end{itemize}

Для гладкой границы и точки $\boldsymbol{\xi}\in\Gamma$ обычно
$$
c(\boldsymbol{\xi}) = \frac{1}{2}.
$$
В угловых точках значение $c(\boldsymbol{\xi})$ определяется внутренним углом границы.

Таким образом, \textbf{объёмная дифференциальная задача заменяется интегральным уравнением по границе}.

Если исходное уравнение содержит правую часть $f$, то в общем случае возникает также объёмный интеграл. Для получения полностью граничной формулировки его преобразуют дополнительными методами (например, методом двойственности или через частное решение) либо используют специальное представление объёмной нагрузки.

\subsection{Дискретизация границы}

Граница $\Gamma$ разбивается на конечное число участков:
$$
\Gamma = \bigcup_{e=1}^{N}\Gamma_e.
$$

Каждый участок называется \textbf{граничным элементом}.

В двумерной задаче граница может быть разбита на отрезки, а в трёхмерной --- на треугольники или четырёхугольники.

В отличие от метода конечных элементов, внутренняя область $\Omega$ не разбивается на элементы. Дискретизируется только её граница.

Это позволяет уменьшить размерность дискретной задачи:
$$
3D \rightarrow 2D,
\qquad
2D \rightarrow 1D.
$$

\subsection{Граничные элементы и аппроксимация неизвестных}

На каждом граничном элементе неизвестные функции $u$ и
$$
q = \frac{\partial u}{\partial n}
$$
заменяются приближёнными функциями.

При кусочно-постоянной аппроксимации:
$$
u(\mathbf{x}) \approx u_e,
\qquad
q(\mathbf{x}) \approx q_e,
\qquad
\mathbf{x}\in\Gamma_e.
$$

При линейной аппроксимации используются узловые значения:
$$
u(\mathbf{x})
\approx
N_1(\mathbf{x})\,u_1
+
N_2(\mathbf{x})\,u_2,
$$
где $N_1$ и $N_2$ --- функции формы.

Таким образом, непрерывные граничные функции заменяются конечным набором неизвестных --- узловыми значениями искомой функции и её нормальной производной.

\subsection{Формирование системы линейных уравнений}

После подстановки аппроксимаций в граничное интегральное уравнение и перехода к дискретной форме получаем:
$$
\sum_{e=1}^{N}
\int_{\Gamma_e}(\ldots)\,d\Gamma.
$$

Интегралы вычисляются численно, например с использованием квадратурных формул Гаусса.

В результате формируется система линейных алгебраических уравнений:
$$
\mathbf{H}\,\mathbf{u}
=
\mathbf{G}\,\mathbf{q}
+
\mathbf{b},
$$
где
\begin{itemize}
	\item $\mathbf{u}$ --- вектор граничных значений $u$;
	\item $\mathbf{q}$ --- вектор граничных значений нормальной производной;
	\item $\mathbf{H}$ --- матрица интегралов от $\partial G/\partial n$;
	\item $\mathbf{G}$ --- матрица интегралов от $G$;
	\item $\mathbf{b}$ --- вклад объёмных источников, если они присутствуют.
\end{itemize}

На каждой части границы задано одно из условий: либо $u$ (Дирихле), либо $q$ (Нейман), либо их линейная комбинация (Робен). С учётом граничных условий известные величины переносятся в правую часть, и система приводится к стандартному виду
$$
\mathbf{A}\mathbf{x}=\mathbf{f}.
$$

После решения системы находятся неизвестные значения на границе.

Зная граничные значения, решение $u$ во внутренних точках области можно восстановить с помощью интегрального представления:
$$
u(\boldsymbol{\xi})
=
\int_{\Gamma} G(\mathbf{x},\boldsymbol{\xi})\,
\frac{\partial u(\mathbf{x})}{\partial n}
\,d\Gamma
-
\int_{\Gamma} u(\mathbf{x})\,
\frac{\partial G(\mathbf{x},\boldsymbol{\xi})}{\partial n}
\,d\Gamma,
\qquad \boldsymbol{\xi}\in\Omega.
$$

\subsection{Сингулярные интегралы}

Одной из особенностей МГЭ является наличие \textbf{сингулярных интегралов}.

Фундаментальное решение $G$ и его производные имеют сингулярность, когда точка наблюдения совпадает с точкой источника:
$$
\mathbf{x}\rightarrow\boldsymbol{\xi}.
$$

Поэтому некоторые граничные интегралы становятся несобственными или сингулярными. При этом диагональные элементы матрицы $\mathbf{H}$ содержат коэффициент $c(\boldsymbol{\xi})$ (для гладкой границы $c=1/2$).

Для их вычисления применяют:
\begin{itemize}
	\item аналитическое вычисление интегралов;
	\item специальные квадратурные формулы;
	\item регуляризацию;
	\item вычисление главного значения интеграла;
	\item специальные методы численного интегрирования.
\end{itemize}

Корректный учёт сингулярностей является одной из наиболее важных технических задач при реализации МГЭ.

\subsection{Преимущества и ограничения метода}

\textbf{Преимущества:}
\begin{enumerate}
	\item \textbf{Уменьшение размерности задачи.} Дискретизируется только граница:
	$$
	3D\rightarrow2D,\qquad 2D\rightarrow1D.
	$$
	
	\item \textbf{Удобство для бесконечных областей.} Например, при решении задач акустики, механики грунтов и распространения волн.
	
	\item \textbf{Не требуется объёмная сетка.} Это упрощает подготовку модели для задач, в которых внутренняя область имеет сложную геометрию.
	
	\item \textbf{Высокая точность} для гладких границ и достаточно регулярных решений.
\end{enumerate}

\textbf{Ограничения:}
\begin{enumerate}
	\item Необходимо знать фундаментальное решение соответствующего дифференциального оператора.
	
	\item Возникают сложности при наличии объёмных источников и неоднородностей.
	
	\item Необходимо специально обрабатывать сингулярные интегралы.
	
	\item Матрицы МГЭ, как правило, являются \textbf{плотными}, поэтому при большом количестве граничных элементов возрастают требования к памяти и вычислительным ресурсам.
\end{enumerate}

\subsection{Сравнение с методом конечных элементов}

\begin{center}
	\begin{tabular}{|p{0.30\textwidth}|p{0.30\textwidth}|p{0.30\textwidth}|}
		\hline
		\textbf{Характеристика} & \textbf{МГЭ} & \textbf{МКЭ} \\
		\hline
		Где строится сетка & Только на границе & Во всём объёме \\
		\hline
		Размерность сетки & На 1 меньше размерности области & Совпадает с размерностью области \\
		\hline
		Основная формулировка & Интегральная & Вариационная/дифференциальная \\
		\hline
		Бесконечные области & Очень удобен & Требуют специальных приёмов \\
		\hline
		Сложная внутренняя структура & Менее удобен & Очень удобен \\
		\hline
		Сингулярные интегралы & Требуют специальной обработки & Обычно отсутствуют в такой форме \\
		\hline
		Матрица системы & Обычно плотная & Обычно разреженная \\
		\hline
		Фундаментальное решение & Требуется & Не требуется \\
		\hline
		Вычисление внутри области & Через интегральное представление & Непосредственно из решения по элементам \\
		\hline
	\end{tabular}
\end{center}

\textbf{Главное различие:} метод конечных элементов аппроксимирует решение во всей области, тогда как метод граничных элементов сводит задачу к определению неизвестных на её границе.

\textbf{Суть метода:} вместо поиска решения во всём объёме области определяются неизвестные граничные значения, после чего решение во внутренних точках восстанавливается с помощью интегрального представления.

\newpage